% =====================================================================
%  report_skeleton.tex — REVTeX skeleton for the `report` skill (capstone/shareable default).
%
%  COMPILE:  latexmk -pdf report_skeleton.tex
%  REQUIRES (vendored alongside this file; no installs needed):
%     revtex4-2.cls , apsrev4-2.bst , report-macros.sty , refs.bib , figs/
%  NO LOCAL TeX?  ->  zip this folder and upload to Overleaf (full TeX, zero installs).
%
%  Substyle: two-column journal look = `reprint`. For a one-column DRAFT, swap `reprint` -> `preprint`.
%  Quick INTERNAL drafts may instead use a plain `article` template (this skeleton is the capstone form).
% =====================================================================
\documentclass[
  aps, prx,            % APS / Physical Review X substyle
  reprint,             % two-column; use `preprint` for a one-column draft
  amsmath, amssymb,    % load AMS the REVTeX-recommended way (as class options)
  nofootinbib,
  superscriptaddress,
]{revtex4-2}

\usepackage{graphicx}
\usepackage{booktabs}          % + REVTeX's own \ruledtabular for journal tables
\usepackage{xcolor}
\usepackage{hyperref}
\usepackage{report-macros}     % vendored: siunitx/physics/braket-free macros

\graphicspath{{figs/}}

\begin{document}

% ---- 1. Title / authors / abstract ----------------------------------------
\title{\textsf{<Title — the claim, concretely>}}
\author{<Author>}
\affiliation{<Affiliation / lab>}
\date{\today}

\begin{abstract}
% Headline result WITH its number + the one or two key evidence deltas. Honest:
% lead with what is actually shown; state the band, not the optimistic edge.
<abstract>
\end{abstract}

\maketitle

% ---- 2. Introduction & objective ------------------------------------------
\section{Introduction and objective}\label{sec:intro}
% The problem; the goal and the decision it drives; what is claimed; position vs the literature
% (what gap this fills, what NOT to claim).

% ---- 3. Problem framing & scope  (from `frame-problem`) --------------------
\section{Problem framing and scope}\label{sec:scope}
% The physics MODEL: which effects are INCLUDED vs EXCLUDED, and WHY (the physics inventory —
% negligible by scale / out of scope / deferred). Key assumptions and the domain of validity.
% The decomposition / approach taken to the problem.

% ---- 4. Physical & model parameters  (master parameter table) -------------
\section{Physical and model parameters}\label{sec:params}
% EVERY parameter that entered the result, with provenance + justification + uncertainty.
\begin{table}[t]
  \caption{\label{tab:params}%
    All parameters entering the result. Provenance: user-given / fitted / first-principles /
    simulation / assumed / literature. Justification is given per parameter in the text below.}
  \begin{ruledtabular}
  \begin{tabular}{l l l l l}
    Symbol & Meaning & Value (unit) & Provenance & Uncertainty \\
    \colrule
    $<>$ & <> & $<>\unit{<>}$ & <fit/first-principles/sim/...> & $<\pm>$ \\
    $<>$ & <> & $<>$          & <user-given>                   & --- \\
  \end{tabular}
  \end{ruledtabular}
\end{table}
% Then a paragraph each for the few parameters that matter most: the physics, the governing
% equation, how the value was pinned (derivation / fit anchors / simulation), and the evidence.

% ---- 5. Methods, pipeline & derivations -----------------------------------
\section{Methods}\label{sec:methods}
% End-to-end pipeline so it can be re-run. Derive what is specific to this work; cite standard
% results rather than re-deriving. Define the metric / figure of merit precisely.

% ---- 6. Verification & validation  (the report-time check + ledger) -------
\section{Verification and validation}\label{sec:verify}
% How correctness was established: the verification plan executed — analytic limits, numerical
% convergence, independent-method parity, and any APPLICATION-SPECIFIC checks. Include the
% VERIFICATION LEDGER: which load-bearing claims/values were re-checked, how, and the result.
% Flag anything resting on a QUESTIONABLE ASSUMPTION or that is PROVISIONAL.

% ---- 7. Results & evidence ------------------------------------------------
\section{Results and evidence}\label{sec:results}
% The evidence ledger (each claim -> its delta and source); the headline-result figure; what each
% lever bought; uncertainty bands on every reported quantity.

% ---- 8. (optional) Promising alternatives ---------------------------------
% \section{Promising alternatives}\label{sec:alts}  % brief: the road not taken and why.

% ---- 9. Reproducibility ----------------------------------------------------
\section{Reproducibility}\label{sec:repro}
% Exact configs, commands, seeds, data provenance, environment, and file pointers to regenerate
% the headline result.

% ---- 10. Limitations & open questions -------------------------------------
\section{Limitations and open questions}\label{sec:limits}
% Honest accounting: the known ceiling, unresolved disputes, and every provisional /
% assumption-dependent claim marked as such.

% ---- 11. References (verified by cite-check) -------------------------------
\bibliographystyle{apsrev4-2}
\bibliography{refs}   % refs.bib — every entry confirmed against the real source

\end{document}
